We first establish the full-location rate.  On the trimmed straight interface, each arm contains a fixed half-ball about every target point for all sufficiently small bandwidths.  The density bounds make every scaled bivariate polynomial Gram matrix uniformly nonsingular.  Taylor expansion of a \(q\)-H\"older regression to order \(q-1\) gives bias at most \(Ch^q\).  Repeating the finite-mesh conditional sub-Gaussian argument from the rectifiable theorem, now with bivariate polynomial weights of fixed degree, gives stochastic error
\[
 C\sqrt{\frac{\log n}{nh^2}}.
\]
Thus a full-location local-polynomial estimator with
\(h\asymp(\log n/n)^{1/(2q+2)}\) has risk at most \(Cs_{n,q}\).

For the matching lower bound, take the uniform score density and Gaussian-noised Bernoulli outcomes used in the packing proof, but perturb one treatment regression by one of \(M\asymp h^{-1}\) disjoint two-dimensional bumps of width \(h\) and height \(h^q\), centered along the trimmed interface.  Every law remains in the straight class.  For two such laws, the one-observation relative entropy is bounded by a constant times
\(h^{2q}h^2\), so the product relative entropy is at most \(Cnh^{2q+2}\).  With the displayed bandwidth this is at most a sufficiently small multiple of \(\log M\), after reducing the bump amplitude.  The elementary entropy proof of Fano's inequality is as follows: if \(J\) is uniform on the \(M\) alternatives and \(Z\) is the sample, then
\[
 I(J;Z)\leq M^{-1}\sum_{j=1}^M\operatorname{KL}(P_j^{\otimes n},P_0^{\otimes n})
 \leq\alpha\log M,
\]
whereas any decoder with error probability \(p_e\) satisfies
\(I(J;Z)\geq(1-p_e)\log M-\log2\).  Hence \(p_e\) is bounded away from zero.  Sup-norm accuracy smaller than half the bump height would decode \(J\), proving the lower bound \(cs_{n,q}\).  This lower bound applies even to the full-location class, and therefore also to the distance-only class.

Suppose now \(L>4\).  The local-constant upper argument gives the distance upper bound \(Ca_n\).  The canonical-square version of the narrowed angular packing is a subset of \(\mathcal P_{\mathrm{straight}}\), and its exact-recovery argument gives the lower bound \(ca_n\).

Finally suppose \(L=4\).  Since \(\mathcal X_P=[-1,1]^2\) has area four and \(f_P\geq1/4\),
\[
 1=\int_{[-1,1]^2}f_P(z)\,dz\geq1.
\]
Equality forces \(f_P=1/4\) almost everywhere, and continuity makes the equality pointwise.  For a trimmed boundary point and sufficiently small \(s\), each arm-specific half-circle lies inside the square, so \(Q_{t,P,x}(s)=\pi s/4\).  Taylor's formula for \(\mu_{t,P}\), integrated over that half-circle, gives coefficients \(b_{t,k,P}(x)\) such that
\[
 \left|g^{\mathrm{dist}}_{t,P,x}(s)-\mu_{t,P}(x)-\sum_{k=1}^{q-1}b_{t,k,P}(x)s^k\right|\leq Cs^q.
\]
A degree-\((q-1)\) one-sided local polynomial in signed radius therefore has bias \(Ch^q\), denominator order \(nh^2\), and uniform stochastic term \(C\sqrt{\log n/(nh^2)}\).  Balancing gives the distance upper bound \(Cs_{n,q}\).  The full-location lower bound already proved transfers to the smaller distance decision class, so the distance lower bound is \(cs_{n,q}\).  This proves every asserted case and shows exactly why the frozen all-\(L\) distance claim required correction.